\subsection{Image deblurring}
\label{subsubsec:image_deblurring}
We use image deblurring as a compact, reproducible test case. We adopt as initial steps the procedure described in Section~\ref{sec:common_methodology}. Given a blurred and noisy image $x\in\mathbb{R}^n$ (vectorised and scaled to $[0,1]$), the decision variable $z\in\mathbb{R}^n$ represents the recovered image. We adopt the standard convex model
\begin{equation}
\label{eq:deblur_model}
\begin{aligned}
\underset{z\in\mathbb{R}^n}{\text{minimise}}\quad 
& \tfrac{1}{2}\,\|Az-x\|_2^2 \;+\; \rho\,\|z\|_1 \\
\text{subject to}\quad & 0 \le z \le 1 ,
\end{aligned}
\end{equation}
where $A\in\mathbb{R}^{n\times n}$ is a (known) Gaussian–blur operator realised by a two–dimensional convolution, and $\rho>0$ balances data fidelity and sparsity. The box constraint enforces valid pixel intensities.

\emph{First–order method.} We solve \eqref{eq:deblur_model} with a projected proximal–gradient iteration (ISTA with box projection). Writing $f(z)\defeq\tfrac{1}{2}\|Az-x\|_2^2$ and $g(z)\defeq \rho\|z\|_1+\iota_{[0,1]^n}(z)$, the update is
\[
z_{k+1}
= \Pi_{[0,1]^n}\!\Big(S_{\alpha\rho}\big(z_k-\alpha\nabla f(z_k)\big)\Big),
\qquad 
\nabla f(z)=A^\top(Az-x),
\]
with soft–thresholding $S_{\tau}(u)_i=\mathrm{sign}(u_i)\max\{|u_i|-\tau,0\}$ and projection $\Pi_{[0,1]^n}$ implemented by elementwise clipping. We use a cold start $z_0=0$ and a fixed stepsize $\alpha\in(0,1/L)$, where $L=\|A\|_2^2$; for a convolutional $A$, an accurate $L$ is obtained from the maximum squared magnitude of the kernel’s discrete Fourier transform.

\emph{Performance metric.} For calibration we report a single scalar per run as required by the framework: the normalised mean–squared error (NMSE) in decibels between the iterate $z_K$ and the clean image $y^\star$ that generated $x$,
\[
\phi(z_K,x)\;=\;10\log_{10}\!\left(\frac{\|z_K-y^\star\|_2^2}{\|y^\star\|_2^2}\right)\;\; \text{(lower is better)}.
\]
PSNR can be used interchangeably; we use NMSE(dB) to remain consistent with related work.

\emph{Numerical set–up.} Our sample space $\mathcal{X}$ is generated from EMNIST letters ($28\times 28$ greyscale). For each clean image $y^\star\!\in[0,1]^n$ we form
\[
x \;=\; A y^\star + \eta,\qquad 
A \equiv \text{conv2D}(\,\cdot\,, g_{\sigma_b}),\quad 
\eta\sim\mathcal{N}(0,\sigma_n^2 I),
\]
where $g_{\sigma_b}$ is a Gaussian kernel with support size $8$ and standard deviation chosen so that $95\%$ of its mass lies within the support (we use $\sigma_b=1.6$ for $28\times 28$), $\sigma_n=10^{-3}$, and $\rho=10^{-4}$. The operator $A$ and its adjoint $A^\top$ are evaluated by FFT–based circular convolution for speed and deterministic adjointness.

\emph{Implementation notes (Python).} This example is intended to be run directly within the two–step driver.
\begin{itemize}[leftmargin=*, itemsep=1pt]
\item \texttt{operators.py}: \texttt{class BlurOperator}\\
\hspace*{1.2em}\texttt{__init__(kernel):} precompute FFT of $g_{\sigma_b}$ and its flip;\\
\hspace*{1.2em}\texttt{matvec(z)} $\to Az$ and \texttt{rmatvec(u)} $\to A^\top u$ via FFTs.
\item \texttt{methods.py}: \texttt{projected\_ista(A, x, rho, K, alpha)}\\
\hspace*{1.2em}returns $\{z_k\}_{k=0}^K$ with $z_{k+1}=\Pi_{[0,1]^n}\!\big(S_{\alpha\rho}(z_k-\alpha A^\top(Az_k-x))\big)$.
\item \texttt{metrics.py}: \texttt{nmse\_db(z, y\_true)} and \texttt{psnr(z, y\_true)}.
\item \texttt{data.py}: EMNIST loader returning pairs $(y^\star,x)$ with controlled $(\sigma_b,\sigma_n)$.
\item \texttt{experiment\_deblur.py}:\\
\hspace*{1.2em}\texttt{def run(N, K, rho=1e-4, alpha=None, seed=0):}\\
\hspace*{2.4em}sample $N$ images, build $A$, construct $x$, run ISTA for $K{+}1$ iterates from $z_0\!=\!0$, and return $\{\phi_i\}_{i=1}^N$ with $\phi_i=\mathrm{NMSE}_{\mathrm{dB}}(z_K^{(i)},y^{\star (i)})$.
\end{itemize}
Recommended defaults are $K\in[50,200]$ and $\alpha=0.99/L$ with $L$ estimated once from the kernel FFT. This set–up yields a lightweight, fully deterministic benchmark that cleanly instantiates the two–step pipeline without introducing extraneous notation or solver details.